\section{Conclusion}
\label{sec:conclusion}

This paper argues for a concrete shift in how uncertainty in medical image
segmentation is conceived and evaluated: from estimation to decision.
Under this view, an uncertainty map is not an output---it is an input to a
policy, and its quality is meaningful only with respect to that policy.
Three empirical findings anchor the shift.

\textbf{First, the uncertainty source matters far more for decisions than
for segmentation.} TTA and MC Dropout produce near-identical Dice
(0.768 vs.\ 0.763), yet TTA's uncertainty identifies errors far more
reliably (Unc-AUROC 0.881 vs.\ 0.722). Two methods that look equivalent
by segmentation metrics diverge sharply once a decision layer consumes
their output.

\textbf{Second, how uncertainty is consumed matters as much as how it is
generated.} TTA with adaptive deferral removes roughly 80\% of
segmentation error at 25\% deferral; confidence-aware deferral removes
55\% at just 12\%. The decision rule is not an afterthought---it is half
the system.

\textbf{Third, calibration is not a prerequisite for good deferral.}
Temperature scaling does not improve deferral operating points and
worsens ECE. Calibration optimizes a distributional property; deferral
requires a discriminative one. These objectives decouple by construction.

The broader claim is conceptual: \emph{uncertainty is not a property of a
model, but of a decision process.} Two models with identical calibration
can yield decisions of very different quality, and two decision rules can
extract wildly different value from the same uncertainty map. Evaluating
uncertainty in isolation---the default in both Bayesian deep learning and
medical imaging benchmarks---hides this interaction and rewards the wrong
methods. \textbf{Uncertainty without decisions is evaluation without
consequence.}

Future work should extend this analysis to multi-class segmentation and
3D imaging, where deferral at the voxel, slice, or volume level
introduces additional design choices, and to human-in-the-loop evaluation,
where clinicians interact with deferral maps and diagnostic accuracy is
measured directly. The largest open question is whether the decision-aware
view scales from pixels to patient outcomes.
